\section{The Unquenchable Fire and the Disagreement of the Doctors}

In the third chapter of the Book of Daniel, King Nabuchodonosor orders three young Hebrews cast bound into a furnace heated seven times hotter than usual. When the king peers into the flames, he starts up in astonishment: the men walk loose and unharmed in the midst of the fire, ``and not a hair of their head had been singed, nor their garments altered, nor the smell of the fire had passed on them'' (Dan 3:91--94). St.\ Gregory the Great meditates on that miracle in the ninth book of his \textit{Moralia on Job} as a key to the mystery of hell's fire. In Babylon, ``by the wonderful dispensation of our Creator, the property of fire, being modified into an opposite power, at the same time never touched their garments, and yet burnt their chains.''\endnote{St.\ Gregory the Great, \textit{Morals on the Book of Job} IX.102 (IX.66 in PL 75), Library of the Fathers translation.} Fire is a creature, not an independent sovereign; in the hand of its Maker it can cool into a wind of dew for the righteous while snapping their fetters, or burn without consuming and without giving light for the punishment of the wicked.

Beside the pain of loss, Scripture and the Church teach a real \textit{poena sensus}, a pain of sense, answering to the sinner's inordinate attachment to creatures. As the Book of Wisdom lays down the principle of divine retribution, ``by what things a man sinneth, by the same also he is tormented'' (Wis 11:17). Among all scriptural designations of this positive torment, none recurs with greater insistence than fire: the ``unquenchable fire'' that burns the chaff (Matt 3:12), the ``furnace of fire'' (Matt 13:42, 50), the ``everlasting fire, which was prepared for the devil and his angels'' (Matt 25:41), the fire that ``is not extinguished'' (Mark 9:43, 45, 47), the flame in which the rich man cries out in torment (Luke 16:24), and the ``pool burning with fire and brimstone'' of the Apocalypse (Apoc 20:9--10; 21:8). The \textit{Roman Catechism} notes that the words ``into everlasting fire'' express the pain of sense, ``because, like other corporal punishments, amongst which, no doubt, fire produces the most intense pain, it is felt through the organs of sense.''\endnote{\textit{Catechism of the Council of Trent}, Part I, Creed Article VII, trans.\ Donovan, pp.\ 64--65.}

When Catholic theology asks what \textit{kind} of fire this is, it uncovers both a deep consensus and a genuine disagreement among revered Doctors of East and West. The consensus is unmistakable: the fire of hell is not an empty figure of speech meant to disguise the absence of real punishment, nor is it merely the natural discomfort of a bad memory. It signifies a true, objective, divinely ordered affliction that reaches the whole human person---soul immediately after death, and reunited body and soul after the general resurrection. On the precise physical nature of that fire, however, the Latin tradition and a major Eastern Doctor do not speak with a single voice.

In the Latin West, St.\ Augustine laid down the line that shaped thirteen centuries of Western theology. In \textit{The City of God} XXI.9--10, Augustine examines three interpretations held by Christians of his own day: some refer both the fire and the worm solely to spiritual anguish of the soul; others refer the fire literally to the body and the worm figuratively to remorse of soul; and a third view refers both literally to the body. Augustine rejects the first opinion outright because it leaves the risen body without bodily punishment, and after granting readers freedom regarding the worm, he speaks with complete firmness about the fire: ``only, this we believe about those future bodies, that they shall be such as shall certainly be pained by the fire.''\endnote{St.\ Augustine, \textit{The City of God} XXI.9, trans.\ Dods.}

In the next chapter, Augustine faces the sharpest philosophical difficulty raised by that claim: if the fire of Gehenna is material and corporeal, how can it torment the devil and his angels, who are incorporeal spirits? After noting and setting aside the older speculation that demons possess bodies of dense air, Augustine gives an answer rooted in the mystery of human nature itself. If an immaterial human soul can be joined to a material body in this life so intimately that the soul genuinely suffers when the flesh is burned, why cannot Almighty God bind wicked spirits to material fire so that they genuinely suffer from it?

\begin{quote}
Therefore, though the devils have no bodies, yet their spirits, that is, the devils themselves, shall be brought into thorough contact with the material fires, to be tormented by them; not that the fires themselves with which they are brought into contact shall be animated by their connection with these spirits, and become animals composed of body and spirit, but, as I said, this junction will be effected in a wonderful and ineffable way, so that they shall receive pain from the fires, but give no life to them. \dots\ But that hell, which also is called a lake of fire and brimstone, will be material fire, and will torment the bodies of the damned, whether men or devils\dots\ One fire certainly shall be the lot of both, for thus the truth has declared.\endnote{St.\ Augustine, \textit{The City of God} XXI.10, trans.\ Dods.}
\end{quote}

Augustine is notably cautious, in the same chapter, about the flame of the rich man in Luke 16:24 before the resurrection, observing that someone might reasonably take that intermediate flame to be immaterial like the eyes, tongue, and finger in a disembodied vision; yet of the Gehenna pronounced at the final judgment he has no doubt that the fire is material. St.\ Gregory the Great agrees without hesitation in his \textit{Dialogues}: ``I doubt not that the fire of hell is corporeal, since it is certain that bodies are tortured there.''\endnote{St.\ Gregory the Great, \textit{Dialogues} IV.29, as cited in \textit{Summa Theologiae}, Supplement, q.\ 97, a.\ 5, \textit{sed contra}.}

High scholasticism developed Augustine's and Gregory's position with austere precision. In the \textit{Supplement} to the \textit{Tertia Pars} of the \textit{Summa Theologiae}---compiled after St.\ Thomas Aquinas's death from his earlier \textit{Commentary on the Sentences}---Question 97 examines the punishment of the damned article by article. Rejecting the opinion of Avicenna that hell's torments are purely intellectual grief or dreamlike phantasms, Article 5 concludes that ``whatever we may say of the fire that torments the separated souls, we must admit that the fire which will torment the bodies of the damned after the resurrection is corporeal, since one cannot fittingly apply a punishment to a body unless that punishment itself be bodily.''\endnote{\textit{Summa Theologiae}, Supplement, q.\ 97, a.\ 5, trans.\ Fathers of the English Dominican Province (London: Burns, Oates \& Washbourne, 1922).} How can one corporeal fire afflict different sinners in different degrees according to their guilt, when an ordinary hearth burns wood and saint alike by mere natural heat? Because, the reply explains, hell's fire acts as ``the instrument of Divine justice inflicting punishment,'' and an instrument acts by the power and direction of its principal agent---``even so the fire of the furnace is regulated by the forethought of the smith, according as the effect of his art requires.''\endnote{\textit{Summa Theologiae}, Supplement, q.\ 97, a.\ 5, ad 3.}

In Article 6, the \textit{Supplement} presses the question one step further: is the fire of hell of the same species as terrestrial fire? The answer distinguishes between the formal nature of fire and the matter in which fire subsists. ``Accordingly it is clear that the fire of hell is of the same species as the fire we have, so far as the nature of fire is concerned. But whether that fire subsists in its proper matter, or if it subsists in a strange matter, what that matter may be, we know not.''\endnote{\textit{Summa Theologiae}, Supplement, q.\ 97, a.\ 6.} Unlike earthly fire, it requires neither human kindling nor fuel to sustain it, for God Himself sustains its nature---the sense in which the \textit{Supplement} reads Isaias 30:33, ``the breath of the Lord as a torrent of brimstone kindling it''---and, like the risen bodies of the damned themselves, it is preserved incorruptible by the ordinance of divine justice.\endnote{\textit{Summa Theologiae}, Supplement, q.\ 97, a.\ 6, ad 2--3.}

Yet the Catholic tradition does not speak here with absolute uniformity. At the very close of his \textit{Exposition of the Orthodox Faith} (Book IV, Chapter 27), St.\ John of Damascus---the last of the Greek Fathers and a Doctor of the universal Church---summarizes the resurrection and final judgment in words that deliberately set hell's fire apart from earthly combustion:

\begin{quote}
We shall therefore rise again, our souls being once more united with our bodies, now made incorruptible and having put off corruption, and we shall stand beside the awful judgment-seat of Christ: and the devil and his demons and the man that is his, that is the Antichrist and the impious and the sinful, will be given over to everlasting fire: not material fire like our fire, but such fire as God would know.\endnote{St.\ John of Damascus, \textit{Exposition of the Orthodox Faith} IV.27, trans.\ S.\ D.\ F.\ Salmond, NPNF2 vol.\ 9.}
\end{quote}

Damascene's sentence is marked by deep reserve: ``not material fire like our fire, but such fire as God would know.'' He affirms the resurrection of the flesh, the awful tribunal of Christ, and the reality of everlasting fire for both demons and wicked men, while refusing to equate that fire with the material element known on earth. The compiler of the Thomistic \textit{Supplement} faces Damascene's text directly in Question 97, Article 5, Objection 1, and reads it with respect, arguing that Damascene denies not materiality altogether but materiality ``as our fire,'' because hell's fire differs in its properties and acts punitively without corrupting the body's matter.\endnote{\textit{Summa Theologiae}, Supplement, q.\ 97, a.\ 5, ad 1.} Even with that harmonization, a real theological difference remains visible between the Latin school's affirmation of a corporeal fire of the same species as ours and Damascene's more apophatic Greek reserve.

The Church has left that scholastic question open, binding her faithful neither to a literalist physics of subterranean combustion nor to a purely psychological metaphor that dissolves objective punishment. The \textit{Catechism of the Catholic Church} retains our Lord's exact phrase ``eternal fire'' in quotation marks while stating firmly the reality of ``the punishments of hell'' (CCC 1035), and St.\ John Paul II explains that Scripture's symbolic language indicates ``the state of those who freely and definitively separate themselves from God, the source of all life and joy.''\endnote{\textit{Catechism of the Catholic Church}, \S\S 1034--1035; Pope St.\ John Paul II, General Audience, 28 July 1999.} What stands beyond all dispute is the warning that both Augustine and Damascene unite to enforce: whatever the exact mode of its action, the fire prepared for the devil and his angels is real, severe, proportioned to guilt, and utterly unlike anything a sane soul would risk enduring. As St.\ Teresa of Avila wrote after her own mystical glimpse of hell, comparing painted depictions with what she underwent: ``In short, the one is a reality, the other a picture; and all burning here in this life is as nothing in comparison with the fire that is there.''\endnote{St.\ Teresa of Avila, \textit{The Life of St.\ Teresa of Jesus} ch.\ 32, \S 6, trans.\ David Lewis, rev.\ Benedict Zimmerman (London: Thomas Baker, 1904), p.\ 288.}
